\documentclass[aspectratio=169,11pt]{beamer}

\usepackage[spanish,es-nodecimaldot]{babel}
\usepackage{fontspec}
\usepackage{tikz}
\usetikzlibrary{arrows.meta,positioning,shapes.geometric}
\usepackage{xcolor}
\usepackage{listings}
\usepackage{booktabs}

\definecolor{azul}{HTML}{174A7E}
\definecolor{azulclaro}{HTML}{DCECF8}
\definecolor{naranja}{HTML}{D95F02}
\definecolor{verde}{HTML}{2E7D32}
\definecolor{verdeclaro}{HTML}{E8F5E9}
\definecolor{rojo}{HTML}{B33A3A}
\definecolor{fondo}{HTML}{F7F8FA}
\definecolor{obstaculo}{HTML}{343A40}

\usetheme{Madrid}
\setbeamercolor{structure}{fg=azul}
\setbeamercolor{frametitle}{fg=azul}
\setbeamercolor{normal text}{fg=black!80}
\setbeamercolor{block title}{bg=azul,fg=white}
\setbeamercolor{block body}{bg=azulclaro,fg=black!80}
\setbeamertemplate{navigation symbols}{}
\setbeamertemplate{footline}{\hfill\insertframenumber\hspace{0.4cm}\vspace{0.2cm}}

\lstdefinestyle{python}{
  language=Python,
  basicstyle=\ttfamily\small,
  backgroundcolor=\color{fondo},
  frame=single,
  keepspaces=true,
  showstringspaces=false,
  breaklines=true,
  keywordstyle=\color{azul}\bfseries,
  stringstyle=\color{naranja}
}

\title{Laboratorio 7}
\subtitle{Red neuronal pequena y evolucion por algoritmo genetico}
\author{Inteligencia Artificial}
\date{}

\begin{document}

\begin{frame}
  \titlepage
\end{frame}

\begin{frame}{La historia del robot}
  \begin{center}
  \begin{tikzpicture}[node distance=0.3cm]
    \node[draw,rounded corners,fill=azulclaro,minimum width=2.55cm,minimum height=1.9cm,align=center] (uno) {\textbf{Clases 1--2}\\Ruta y fitness};
    \node[draw,rounded corners,fill=orange!15,minimum width=2.55cm,minimum height=1.9cm,align=center,right=of uno] (dos) {\textbf{Clase 3}\\Obstaculos y cruce};
    \node[draw,rounded corners,fill=verdeclaro,minimum width=2.55cm,minimum height=1.9cm,align=center,right=of dos] (tres) {\textbf{Clase 4}\\Percepcion local};
    \node[draw,rounded corners,fill=orange!25,minimum width=2.55cm,minimum height=1.9cm,align=center,right=of tres] (cuatro) {\textbf{Clase 7}\\Red neuronal};
    \draw[->,very thick,azul] (uno) -- (dos);
    \draw[->,very thick,azul] (dos) -- (tres);
    \draw[->,very thick,azul] (tres) -- (cuatro);
  \end{tikzpicture}
  \end{center}
  \vspace{0.35cm}
  \begin{block}{Pregunta de hoy}
    Podemos conservar la percepcion local y cambiar la forma en que el robot representa su politica?
  \end{block}
\end{frame}

\begin{frame}{Que conservamos y que cambia}
  \begin{columns}[T]
    \column{0.46\textwidth}
      \begin{block}{Se conserva}
        \begin{itemize}
          \item La cuadricula y los obstaculos.
          \item Los sensores U, D, L y R.
          \item El recorrido del robot.
          \item El fitness y el algoritmo evolutivo.
        \end{itemize}
      \end{block}
    \column{0.46\textwidth}
      \begin{block}{Cambia}
        \begin{itemize}
          \item La tabla de reglas desaparece.
          \item Cada individuo es una red.
          \item Los genes son pesos numericos.
          \item Se usa PyTorch para decidir.
        \end{itemize}
      \end{block}
  \end{columns}
  \vspace{0.25cm}
  \begin{alertblock}{Idea central}
    La evolucion sigue buscando una buena politica; ahora la politica esta parametrizada por una red neuronal.
  \end{alertblock}
\end{frame}

\begin{frame}{La observacion local}
  \begin{columns}[T]
    \column{0.52\textwidth}
      \begin{center}
      \begin{tikzpicture}[x=0.75cm,y=0.75cm]
        \foreach \x in {0,...,4} {\foreach \y in {0,...,4} {\fill[white] (\x,\y) rectangle ++(1,1);\draw[black!15] (\x,\y) rectangle ++(1,1);}}
        \fill[obstaculo] (2,1) rectangle ++(1,1);
        \fill[obstaculo] (1,3) rectangle ++(1,1);
        \fill[azul!25] (2,2) rectangle ++(1,1);
        \node at (2.5,2.5) {robot};
        \draw[->,very thick,rojo] (2.5,2.5) -- (2.5,3.35) node[above] {U};
        \draw[->,very thick,verde] (2.5,2.5) -- (2.5,1.65) node[below] {D};
        \draw[->,very thick,naranja] (2.5,2.5) -- (1.65,2.5) node[left] {L};
        \draw[->,very thick,azul] (2.5,2.5) -- (3.35,2.5) node[right] {R};
      \end{tikzpicture}
      \end{center}
    \column{0.4\textwidth}
      \begin{itemize}
        \item 0: celda libre.
        \item 1: obstaculo o borde.
        \item Entrada: cuatro bits.
        \item Salida: cuatro acciones.
      \end{itemize}
      \begin{block}{Percepcion parcial}
        El robot no ve el tablero completo ni recibe una ruta.
      \end{block}
  \end{columns}
\end{frame}

\begin{frame}{La red neuronal pequena}
  \begin{center}
  \begin{tikzpicture}[node distance=1.2cm]
    \node[draw,rounded corners,fill=azulclaro,inner sep=8pt] (entrada) {4 sensores};
    \node[draw,rounded corners,fill=orange!20,right=of entrada,inner sep=8pt] (oculta) {8 neuronas + ReLU};
    \node[draw,rounded corners,fill=verdeclaro,right=of oculta,inner sep=8pt] (salida) {4 valores de accion};
    \draw[->,very thick,azul] (entrada) -- (oculta);
    \draw[->,very thick,azul] (oculta) -- (salida);
  \end{tikzpicture}
  \end{center}
  \vspace{0.4cm}
  \begin{block}{Decision}
    La red calcula un valor para U, D, L y R. El robot ejecuta la accion con mayor valor mediante \texttt{argmax}.
  \end{block}
  \begin{alertblock}{Representacion}
    En la politica tabular evolucionabamos reglas. Aqui evolucionamos todos los pesos de la red.
  \end{alertblock}
\end{frame}

\begin{frame}[fragile]{El codigo de la politica}
\begin{lstlisting}[style=python]
class RedPercepcion(nn.Module):
    def __init__(self, hidden_size=8):
        super().__init__()
        self.hidden = nn.Linear(4, hidden_size)
        self.out = nn.Linear(hidden_size, 4)

    def forward(self, x):
        return self.out(torch.relu(self.hidden(x)))
\end{lstlisting}
  \begin{block}{Lectura del codigo}
    Cuatro entradas representan la percepcion local; la capa oculta combina la informacion y cuatro salidas compiten por controlar el siguiente movimiento.
  \end{block}
\end{frame}

\begin{frame}{Como evoluciona una red}
  \begin{center}
  \begin{tikzpicture}[node distance=0.65cm]
    \node[draw,rounded corners,fill=azulclaro,inner sep=7pt] (pob) {40 redes};
    \node[draw,rounded corners,fill=orange!15,right=of pob,inner sep=7pt] (eval) {evaluar fitness};
    \node[draw,rounded corners,fill=verdeclaro,right=of eval,inner sep=7pt] (sel) {elitismo y padres};
    \node[draw,rounded corners,fill=orange!20,right=of sel,inner sep=7pt] (var) {cruce y mutacion};
    \draw[->,very thick,azul] (pob) -- (eval);
    \draw[->,very thick,azul] (eval) -- (sel);
    \draw[->,very thick,azul] (sel) -- (var);
    \draw[->,very thick,azul] (var) to[bend left=30] (pob);
  \end{tikzpicture}
  \end{center}
  \vspace{0.35cm}
  \begin{block}{Fitness}
    \centering
    $F = 500 - 35d + 4p - 30c - 3r + 2000m$
  \end{block}
  \begin{center}
    Se premia acercarse y llegar; se penalizan choques y repeticiones.
  \end{center}
\end{frame}

\begin{frame}{Por que guardar el mejor individuo}
  \begin{columns}[T]
    \column{0.45\textwidth}
      \begin{block}{Durante el entrenamiento}
        \begin{itemize}
          \item Se prueban muchas redes.
          \item Se conserva la mejor.
          \item La evolucion puede terminar.
        \end{itemize}
      \end{block}
    \column{0.45\textwidth}
      \begin{block}{Despues del entrenamiento}
        \begin{itemize}
          \item Se guardan los pesos en \texttt{mejor\_red.pt}.
          \item Se carga la misma politica.
          \item Se prueba en otro mapa.
        \end{itemize}
      \end{block}
  \end{columns}
  \vspace{0.25cm}
  \begin{alertblock}{Separar aprender y probar}
    Esto permite medir si la politica funciona solo en el mapa de entrenamiento o si responde bien a una nueva distribucion de obstaculos.
  \end{alertblock}
\end{frame}

\begin{frame}[fragile]{Comandos del laboratorio}
\begin{lstlisting}[style=python]
# Entrenar y guardar
python robot_red_neuronal_evolutiva.py \\
  --generaciones 40 --guardar-mejor mejor_red.pt

# Cargar sin volver a evolucionar
python robot_red_neuronal_evolutiva.py \\
  --cargar-red mejor_red.pt
\end{lstlisting}
  \begin{block}{Mapa alternativo}
    Se modifica manualmente \texttt{OBSTACULOS}, \texttt{INICIO} o \texttt{META} en el archivo y se repite el segundo comando.
  \end{block}
\end{frame}

\begin{frame}{Pregunta experimental}
  \begin{center}
    \Large Una politica que funciona en un mapa, funcionara en otro?
  \end{center}
  \vspace{0.3cm}
  \begin{enumerate}
    \item Entrenar con el mapa original.
    \item Guardar el mejor individuo.
    \item Cambiar solo los obstaculos.
    \item Cargar la red y observar la trayectoria.
    \item Comparar llegada, choques, pasos y ciclos.
  \end{enumerate}
  \begin{alertblock}{Hipotesis}
    Como la red solo ve vecinos, puede reutilizar comportamientos locales, pero no conoce el mapa completo ni la distancia global a la meta.
  \end{alertblock}
\end{frame}

\begin{frame}{Cierre: de reglas a representaciones}
  \begin{itemize}
    \item La clase anterior introdujo una politica que reacciona a la percepcion.
    \item Este laboratorio reemplaza la tabla por una red neuronal pequena.
    \item El algoritmo genetico sigue siendo el mecanismo de aprendizaje.
    \item Guardar el individuo permite separar entrenamiento y prueba.
    \item El mapa nuevo introduce una primera pregunta de generalizacion.
  \end{itemize}
  \vspace{0.35cm}
  \begin{block}{Idea final}
    No cambiamos el problema para ocultar la evolucion: cambiamos la representacion para estudiar que puede aprender una politica parametrica.
  \end{block}
\end{frame}

\end{document}
